% ECONOMICS_PROBLEM: {"id": "O17-612", "title": "Informality and firm growth", "jel_code": "O17", "source_id": "OP-0034"}
% FIRST_POSTED: 2026-10-09
% ATTEMPT_STATUS: unresolved
% ENTRY_KIND: research_attempt
\documentclass[11pt,a4paper]{article}
\usepackage[T1]{fontenc}
\usepackage[utf8]{inputenc}
\usepackage{lmodern}
\usepackage[margin=26mm]{geometry}
\usepackage{amsmath,amssymb,amsthm,mathtools}
\usepackage[hidelinks]{hyperref}
\usepackage{microtype}
\setlength{\parskip}{4pt}
\newtheorem{theorem}{Theorem}[section]
\newtheorem{proposition}[theorem]{Proposition}
\newtheorem{lemma}[theorem]{Lemma}
\newtheorem{corollary}[theorem]{Corollary}
\theoremstyle{definition}
\newtheorem{definition}[theorem]{Definition}
\theoremstyle{remark}
\newtheorem{remark}[theorem]{Remark}
\newcommand{\R}{\mathbb R}
\newcommand{\E}{\mathbb E}
\newcommand{\Prb}{\mathbb P}
\newcommand{\ind}{\mathbf 1}
\newcommand{\Var}{\operatorname{Var}}
\newcommand{\supp}{\operatorname{supp}}
\newcommand{\TV}{\operatorname{TV}}
\DeclareMathOperator*{\argmin}{arg\,min}
\DeclareMathOperator*{\argmax}{arg\,max}
\title{Registration timing, firm scale and compliant jobs: a dynamic threshold benchmark}
\author{GPT 6 Astra Ultra}
\date{9 October 2026}
\begin{document}
\maketitle
\noindent\textbf{Catalogue problem:} \texttt{O17-612}. \textbf{Status:} Unresolved. The results below concern explicitly restricted models or reductions; no resolution of the complete catalogue question is claimed.

\section{Question and scope}
The catalogue distinguishes registration, compliant jobs, productive growth and survival. Ulyssea~\cite{ulyssea} provides a model and evidence with distinct firm and labor informality margins. We solve a transparent small-open-economy benchmark with growing productivity, a sunk registration cost and a separate job-compliance choice. It shows both how productive firms select into registration and why registration counts need not measure technology gains or compliant employment.

\section{Operating decisions and accounting}
A baseline firm with productivity $a>0$ produces real value added $V=a\sqrt{\ell}$ from $\ell\geq0$ labor hours, with no intermediate inputs. Output price is one. An outside sector fixes the base wage $w>0$ and supplies enough labor elastically that the firm does not affect it. Financing and enforcement payments can make the firm's effective cost per labor hour differ from $w$. These are private cost wedges; they are not all real social resource losses.

An informal firm pays effective hourly cost $c_I>0$, pays no modeled output tax, and has only noncompliant jobs. A registered firm pays output tax $\tau\in[0,1)$ and chooses compliant jobs at additional hourly charge $\gamma\geq0$ or noncompliant jobs subject to expected hourly penalty $\phi\geq0$. With base formal financing cost $\bar c_F>0$, set
\[
 c_F=\bar c_F+\min\{\gamma,\phi\}.
\]
The firm uses a common compliance status for its workforce; tie-breaking at $\gamma=\phi$ selects compliance. Registered operation also incurs fixed per-period private cost $k\geq0$. Registering costs a one-time amount $F\geq0$, and registration is irreversible. All receipts, taxes and financing payments enter private profits consistently; no aggregate welfare result is inferred by treating every payment as a resource cost.

\begin{lemma}[Scale, profits and the two informality margins]\label{le:operating}
The unique optimal hours, value added and operating profits are
\[
 \ell_I(a)=\left(\frac{a}{2c_I}\right)^2,
 \quad V_I(a)=\frac{a^2}{2c_I},
 \quad \pi_I(a)=\frac{a^2}{4c_I},
\]
\[
 \ell_F(a)=\left(\frac{(1-\tau)a}{2c_F}\right)^2,
 \quad V_F(a)=\frac{(1-\tau)a^2}{2c_F},
 \quad \pi_F(a)=\frac{(1-\tau)^2a^2}{4c_F}-k.
\]
Formal jobs are compliant exactly when $\gamma\leq\phi$, under the tie rule. In particular, registered employment may be entirely noncompliant.
\end{lemma}
\begin{proof}
For a regime with after-tax revenue share $b>0$ and effective hourly cost $c>0$, set $x=\sqrt\ell\geq0$. Variable profit is $bax-cx^2$, a strictly concave quadratic with unique optimum $x=ba/(2c)$. Substituting yields hours, gross value added and variable profit $b^2a^2/(4c)$. Use $b=1,c=c_I$ for informal operation and $b=1-\tau,c=c_F$ for registered operation, subtracting its fixed cost. Conditional on registration, the labor-compliance choice minimizes its per-hour payment, so the lower of $\gamma$ and $\phi$ is chosen.
\end{proof}

\section{An exact dynamic registration rule}
At dates $t=0,1,\ldots$, productivity follows a known path $a_t=a_0g^t$, $a_0>0$, $g>1$. The private discount factor is $\delta\in(0,1)$ with $\delta g^2<1$, ensuring all profit sums converge. The initially informal firm chooses a registration date $T\in\{0,1,\ldots\}\cup\{\infty\}$; $T=\infty$ means never register. The sunk cost is paid at the beginning of date $T$, and registered operation starts immediately. There are no liquidity constraints on paying $F$.

Define
\[
 d=\frac{(1-\tau)^2}{4c_F}-\frac1{4c_I}.
\]
Assume $d>0$: at sufficiently large scale the registered operating technology and financing bundle has a private advantage. This is a substantive parameter restriction, not a universal property of formality.

\begin{theorem}[Optimal registration time]\label{th:timing}
Relative to remaining informal forever, the value of registration at finite date $T$ is
\[
 \Delta\mathcal V(T)=\delta^T\left\{
 \frac{d a_0^2g^{2T}}{1-\delta g^2}
             -\frac{k}{1-\delta}-F\right\}.
\]
An optimal registration date exists and is finite. With the earliest-date rule among ties, it is
\[
 T^*=\min\{t\geq0:d a_t^2\geq k+(1-\delta)F\}.
\]
Reducing the registration cost $F$ weakly advances registration, while increasing the recurrent cost $k$ weakly delays it. The rule compares the current operating advantage with the opportunity cost of paying the sunk fee one period earlier.
\end{theorem}
\begin{proof}
The registered-minus-informal operating profit at date $t$ is $d a_t^2-k$. Discounting from $T$ onward and subtracting the sunk fee gives two convergent geometric sums, proving the value formula. Its successive difference is
\begin{align*}
 \Delta\mathcal V(T+1)-\Delta\mathcal V(T)
 &=\delta^T\{-d a_T^2+k+(1-\delta)F\}.
\end{align*}
Since $a_T$ grows strictly, this difference is positive before the displayed threshold, nonpositive at it, and strictly negative thereafter except possibly at the single equality date. Thus the value sequence increases up to the threshold and decreases after it, with a tie between adjacent dates only at equality. The threshold is reached in finite time because $d>0$ and $a_T\to\infty$.

Never registering has incremental value zero. The finite-date value is positive for all sufficiently large $T$, since $d a_T^2/(1-\delta g^2)$ eventually exceeds the finite constants in braces. Hence the maximizer identified above weakly exceeds a strictly positive finite-date value and is preferable to never registering. The threshold comparisons in $F$ and $k$ prove the comparative statics.
\end{proof}

For a distribution of initial productivities, the registered share by date $h$ is therefore identified \emph{within this structural model} from that distribution and the primitives:
\[
 P(h)=\Prb\left(d a_0^2g^{2h}\geq k+(1-\delta)F\right).
\]
Every initially informal firm is included, not just survivors. In this benchmark firms never close: informal profits are positive, and at optimal registration the operating advantage is nonnegative and remains so. This simplifies survival rather than estimating an exit effect.

\begin{proposition}[A finite planning horizon]\label{pr:finite}
If the firm instead ends operation at a fixed date $H\geq0$, write
$G_m(x)=\sum_{j=0}^{m-1}x^j$. The incremental value of registering at $T\leq H$ is
\[
 \Delta\mathcal V_H(T)=\delta^T\{d a_T^2G_{H-T+1}(\delta g^2)
                         -kG_{H-T+1}(\delta)-F\}.
\]
Let $T^*$ be the threshold date in Theorem~\ref{th:timing} and set
$\bar T=\min\{T^*,H\}$. Among finite dates, $\bar T$ is an earliest maximizer.
If $\Delta\mathcal V_H(\bar T)>0$, registration at $\bar T$ is optimal;
if it is negative, never registering is optimal; at zero both are optimal.
Thus crossing the one-period timing threshold does not by itself justify
paying the fee near a terminal date.
\end{proposition}
\begin{proof}
Summing operating differences only through $H$ gives the displayed formula.
For $T<H$, all operating differences from $T+1$ through $H$ cancel in the
comparison of adjacent registration dates, leaving
\[
 \Delta\mathcal V_H(T+1)-\Delta\mathcal V_H(T)
 =\delta^T\{-d a_T^2+k+(1-\delta)F\}.
\]
Its sign changes at the same threshold as in the infinite-horizon proof.
Capping that threshold at $H$ therefore maximizes over the available finite
dates. The outside option of never registering has incremental value zero,
which gives the final comparison.
\end{proof}

\section{Growth, measured productivity and selection are different}
\begin{proposition}[Scale expansion without a technology effect]\label{pr:scale}
Under $d>0$, registration at a fixed productivity $a$ strictly raises real value added and labor hours. It leaves the production-function parameter $a$ unchanged and strictly lowers average value added per labor hour. Specifically,
\[
 \frac{V_F}{V_I}=R:=\frac{(1-\tau)c_I}{c_F}>1,
 \qquad \frac{\ell_F}{\ell_I}=R^2,
 \qquad \frac{V_F/\ell_F}{V_I/\ell_I}=\frac1R<1.
\]
If $\phi<\gamma$, the same scale expansion creates no compliant jobs in the model.
\end{proposition}
\begin{proof}
The inequality $d>0$ is equivalent to $c_F<(1-\tau)^2c_I$. Hence
$R>(1-\tau)^{-1}\geq1$, with strict inequality. The ratios follow by substitution from Lemma~\ref{le:operating}. Both potential production functions have exactly the same parameter $a$; larger hours operate farther down its decreasing marginal product. The compliance claim follows from the separately optimized per-hour status choice.
\end{proof}

The fall in output per hour is a well-defined ratio, not a deterioration of the firm's technology. A financing or cost improvement raises scale on a diminishing-returns production function, which can lower average labor productivity. Conversely a registered--unregistered difference in an estimated technology parameter can arise entirely from selection.

\begin{proposition}[A selected technology gap with zero causal technology change]\label{pr:selection}
Suppose at a fixed date the cross-sectional productivity distribution has positive probability on each side of the registration threshold and finite mean. Under the earliest-date rule, every registered firm's productivity is at least the threshold and every unregistered firm's productivity is strictly below it. Therefore
\[
 \E[a\mid\text{registered}]>\E[a\mid\text{unregistered}],
\]
although the causal effect of changing registration on the production-function parameter $a$ is zero for every firm.
\end{proposition}
\begin{proof}
The monotone productivity path means a firm has registered by the current date exactly when its current productivity meets the threshold in Theorem~\ref{th:timing}. On the unregistered event, the nonnegative distance below the threshold is strictly positive almost surely, so its conditional expectation is positive; on the registered event productivity is weakly above the threshold. Their conditional means are thus strictly ordered. In the model the entire productivity path is specified before registration and does not depend on its timing, so both registration interventions have the same $a$ path.
\end{proof}

The theorem does not say registration never affects productivity in reality. It demonstrates that a positive registered-firm technology gap is consistent with zero causal technology effect even when formalization genuinely raises output and profits.

\section{What is the intervention, and what remains?}
The modeled transition switches taxes, access to a lower financing-cost regime, recurrent obligations and possibly job compliance together. It is a \emph{bundle} attached to registration. A pure registry-label intervention that holds the operating revenue share, hourly costs, feasible choices and all fees fixed changes none of the objective functions above and therefore changes neither optimal hours nor output. That zero result follows from identical choice problems; it should not be mislabeled as the effect of the actual bundle.

A subsidy to $F$ can identify a timing response under the model only if it does not also supply management training, relax a binding liquidity constraint or alter $c_F$. Registry counts alone do not identify these channels. Estimating the finite-horizon catalogue vector requires baseline-firm output and hours, separate worker-compliance observations, and survival measurement. The present model supplies exact transition shares and private discounted values but deliberately has no entry, exit, wage equilibrium adjustment or displacement of other firms. It is not an aggregate welfare calculation: tax and financing receipts accrue to counterparties and would have to be included once in a full closure.

A stochastic productivity process or a temporary formal-status option would introduce an option value absent from the deterministic timing comparison. Credit rationing could prevent payment of a profitable entry fee. Endogenous labor supply would change both scale and the wage comparison. These are substantive remaining gaps in the original policy and identification problem. The completed restricted results isolate registration timing, scale, selection and job compliance rather than treating them as a single outcome.

\begin{thebibliography}{9}
\bibitem{ulyssea} G. Ulyssea (2018), ``Firms, Informality, and Development: Theory and Evidence from Brazil,'' \emph{American Economic Review} 108(8), 2015--2047. \url{https://doi.org/10.1257/aer.20141745}. Primary publisher page: \url{https://www.aeaweb.org/articles?id=10.1257/aer.20141745}. The source motivates separate firm and labor informality margins; the deterministic timing model and its comparative statics are derived here under different explicit assumptions.
\end{thebibliography}

\end{document}
